\subsection*{任意键先按项确认}

前一节直接按完整 S 工作。活动 E1 已受理 n 个任意 K，逐笔确认可以回收各项大响应，$M_{\min}(S)$ 则继续保存原身份、V/X、真实 P 和退休状态。同 S 异参在原输入处报告冲突，同参从原 P 查询或取得 \code{result\_retired}。只要接口还要准确识别每个任意键，n 项身份依据就随 n 增长，通常记为 $O(n)$。这份依据可以由已有不可变业务记录的可靠身份索引或开放前完成的受控重建承担，未必另加 n 行独立去重表。\footnote{大 O 记号描述输入规模增长时的数量级，此处按保存的记录数计，省略固定常数。键、输入和计数的实际编码长度另计。}

若每项缩减记录平均 m 字节、尚未确认的完整响应有 w 项、每份响应平均 r 字节，会话与配额占 c 字节，则 $L\approx n m+w r+c$。逐笔确认减少 w，会话配额限制 n，关闭 E1 后的整段回收才让入口不再依赖这 n 项依据。

\CTwelveFigure{ch12-space}{逐笔确认回收响应，活动任意键继续保存 n 份身份依据，关闭会话再整段回收}{fig:ch12-space}

如果希望连已确认前缀的逐键已见集合也压缩，接口要提供可核验的顺序，而非给任意键再附一个可随意改写的数字。

\subsection*{连续确认为什么可以压成一个数}

另一条路线把 K 本身规定为登记范围内的单调序号，或规定为由 E1、协议客户端和序号 q 唯一生成的可验证编码。完整 S 的四分量保持，服务可以从 K 自身核对 q，不需要为任意字符串另存反向查找。客户端在首次发送前持久保存原序号和输入，序号不重用。

设三个序号 1、2、3 都已经完成。客户端先持久收到 2、3 的结果，1 的响应却丢失。最大已完成序号为 3，最大已收到序号也是 3，但这两个数字都不能授权服务删除 1 的原结果。客户端下一步要恢复的是 1。

RIFL 用 \code{firstIncomplete} 表示尚未收到响应的最小序号。并发 Raft 客户端的回收说明也采用相同的连续范围~\cite{rifl,ongaro}。\footnote{RIFL 的 RequestTracker 以收到响应维护 \code{firstIncomplete}。本书购买客户端会跨进程恢复任务，因此把可确认条件明确为已持久取得完整终态响应，并将确认解释为主动结束相应回放。服务端独立保留原身份拒绝依据。}
本书把确认加强为客户端已经持久取得结果，记 f 为尚未持久取得终态响应的最小已分配序号，$h_{\rm ack}=f-1$。当已分配的全部序号都有持久结果时，f 取下一个待分配序号。$h_{\rm ack}$ 称为确认水位，表示 1 至 $h_{\rm ack}$ 的结果都已持久取得。它只压缩这一段连续事实。

\begin{BookCode}
已分配序号：     1       2       3       4
服务端状态：     完成    完成    完成    尚未到达
客户端持久结果： 缺失    已有    已有    缺失
firstIncomplete = 1，确认水位 h_ack = 0

序号 1 的原结果补齐并持久保存以后：
firstIncomplete = 4，确认水位 h_ack = 3
\end{BookCode}

\CTwelveFigure{ch12-original06}{已完成最大序号为 3 时，确认水位为 0，补齐 1 后才连续推进到 3}{fig:ch12-prefix}

服务在原受理域保存 $h_{\rm ack}$，并令它单调前进。$h_{\rm ack}$ 以下迟到的网络副本得到已退休响应，业务效果保持。RIFL 将这种来件归入 \code{STALE}，其中既包括客户端不合规则的重发，也包括网络里早已发出、很久以后才到达的副本~\cite{rifl}。客户端确认以后，水位承担了识别这些迟到副本的工作。

调用方确认范围具有明确输入规则：所确认的结果已经持久取得，并同意结束回放。服务能够检查对应序号是否已经形成终态，却无法从一条网络消息证明调用方磁盘确已写入。虚假确认可能使调用方失去原结果。只要服务保留下面的身份拒绝依据，这类错误也不会让同一 S 再次首次执行。服务端的效果上限与客户端自己的结果保存责任分别成立。

对于这种序号键，服务把已经确认的连续前缀 1 至 $h_{\rm ack}$ 标为退休，迟到 q 只需与水位比较。这个已确认前缀的成员资格摘要只占一个整数，按机器字数计可写为 $O(1)$。未确认窗口、乱序空洞、各项输入及结果、会话登记和业务证据分别保存；整数自身的编码位数也会随值增长。如果接口在退休后继续承诺报告每个旧 q 的具体异参冲突，原输入或等价证据也继续保留。仅保留水位的接口统一拒绝退休身份，结束对其具体异参的区分。

\subsection*{保留任意键时，外加序号有什么代价}

有的客户端库固定采用 \code{client\_id/seq}，业务又必须接受任意 K。这时可以选择适配 U：例如 \code{(wallet:E1,1)} 与原 S1 双向唯一绑定。首次分配持久保存，重试使用原编码，服务把 U 与 M 共同仲裁。S1 若被重新编码为 2，反向绑定报告冲突；序号 1 改指 K2，正向绑定报告冲突。两条检查保护服务已经建立的单一绑定。

\CTwelveFigure{ch12-original07}{任意键外加序号的扩展分支保留双向 U，防止旧 K1 重编码为水位以上的新序号}{fig:ch12-reencode}

在这条扩展分支，删除 M(S1) 和 U 两个方向、仅保留水位 1，会让携带旧 K1 的序号 2 再次进入首次分支。水位能回收连续响应，逐键身份和输入成本依旧存在，空间式中 m 还要计入 U。因而任意 K 的主路径直接逐项确认即可。后面为了展开空洞和 SKIP，将同时说明可核验序号键与这条可选适配路径，涉及 U 的行只属于后者。旧任意键会话按原规则结束，新会话若选择序号键，使用明确的新动作版本和键格式。

\subsection*{会话代次怎样结束一整段身份}

当 E1 关闭，E1 下所有尚未来得及首次受理的 K 都失去入口资格。此时完成操作的 $M_{\min}$ 和 U 可以整段回收，因为迟到请求先被 E1 的资格判断拦住。未完成 C/T 继续按原关联保留。已经形成的账务与订单证据，按其业务保存规则另行处理。

有些系统把会话登记的一代称为 epoch，即代次。这里它标记某次登记所拥有的准入范围。E1 和 E2 可以属于同一账户，却是不同的一代。它与 Raft 的领导者任期、T1 的 \code{lease\_epoch} 分别作用于会话准入、复制领导权和恢复任务领取，比较时各用自己的作用对象。

如果每个关闭会话都永久保存一行，逐笔历史可以收缩为逐会话历史，但会话总数继续增长。本书允许进一步删除关闭行，前提是入口对缺失会话一律拒绝，登记接口只分配新的身份。登记器可为账户与动作保存持久递增的分配计数，分配下一代并持久提交后才签发许可。已经发出的代次永不重新分配，计数随副本与快照保存。

\CTwelveFigure{ch12-original08}{E1 关闭后逐键依据可以整段回收，缺失 E1 继续拒绝，登记器只签发新的 E2}{fig:ch12-epoch}

活动登记表保存当前有效会话，计数保护身份分配的历史。许可即使能验证真伪，也只证明曾经签发，入口还要找到当前活动登记并核对范围与截止。持有旧许可的请求不能要求登记器补建 E1。Ongaro 也明确要求命令引用缺失会话时返回错误，后面将用一条重连历史说明这项决定~\cite{ongaro}。

计数也有容量与持久性成本。固定宽度计数到达上限时，登记停止，或通过明确迁移进入新的、不会与旧范围碰撞的命名空间。删除计数、从旧备份恢复较小值，或者按进程重启重新从 1 分配，都会破坏从未重用的前提。本章以登记历史持续保留为条件，第十四章将处理这种历史回退。

\subsection*{固定截止怎样进入准入次序}

本书 E1 采用固定截止和不可撤回关闭。RIFL 则允许客户端续租，其实现只持久保存租约存在，续租截止留在内存，租约服务重启后会续上已有租约~\cite{rifl}。这能服务于它的活跃客户端模型。本书将 E1 的固定截止持久保存，重启恢复原截止和关闭决定。

先看数据库路径。A 与会话关闭都锁定同一登记行。A 在锁内核对资格和权威时间，取得准入决定后保存 P/C/T。关闭事务取得同一锁，将 E1 置为关闭。两者按这个锁形成先后：A 已经取得合法准入并持锁时，关闭等待它提交或回滚；关闭先提交，后来的 A 读到关闭并拒绝。每一次 A 独立检查固定截止，关闭任务迟到时也按相同时间边界判断。

时间的来源要给出保证。若固定截止使用 UTC，准入服务可以从受信时间源取得当前时刻区间 $[t_{\mathrm{low}},t_{\mathrm{high}}]$，实际时刻位于该区间。只有 $t_{\mathrm{high}}<t_{\rm deadline}$ 时，才确认此刻严格早于截止 $t_{\rm deadline}$。时间误差跨越 $t_{\rm deadline}$ 时，入口停止新的受理。$t_{\mathrm{low}}\ge t_{\rm deadline}$ 时，可以明确返回过期。时间源不可用或时钟回拨而无法建立区间时，入口暂停准入并报告暂时不可用，直到恢复可信时间。

\CTwelveFigure{ch12-time-interval}{三个可信时间区间对应受理、暂缓和过期，判断使用同一固定截止 $t_{\rm deadline}$}{fig:ch12-time-interval}

下面的事务展示会话与首次 A 的仲裁，内部已有身份恢复、报价、P/C/T 共同保存沿用前章。\code{checked\_upper} 是锁内取得的受信时间上界，客户端传入的时间戳不能填入它：

\begin{BookCode}
BEGIN
锁定原 E1 登记行
若行缺失或状态为 closed：回滚并返回 operation_expired
核对账户、动作、许可和登记的固定 t_deadline
在锁内取得可信 [checked_lower, checked_upper]
若时间依据不可用：回滚并返回 unavailable
若 checked_lower >= t_deadline：回滚并返回 operation_expired
若 checked_upper >= t_deadline：回滚并返回 cutoff_uncertain
按完整 S 和可选 U 恢复或仲裁
若首次受理成功：共同写入 P/C/T、原输入和固定映射
若原终态或冲突：按原记录决定本次响应
COMMIT
提交确认后发送实际保存或恢复所得的响应
\end{BookCode}

资格判断时刻定义为持锁取得准入决定的时刻。A 可以在截止前取得决定，随后提交稍晚。它在整个事务中持锁，关闭等待这项事务结束。若业务要求连提交完成也必须早于 UTC 的 $t_{\rm deadline}$，就要增加能够在提交边界执行的截止机制，或为事务耗时保留已证明的上界。

这里沿用第~\ref{section:08-2}~节先取行锁、再读可信时间区间的规则。新增对象是 E1 的准入登记行：等待锁的时间不能计作已经取得的资格。PostgreSQL 18.6 的 READ COMMITTED 路径在行锁返回后取得判期依据，事务开始时间不承担这项职责。

\CTwelveFigure{ch12-original09}{同一登记锁排序准入与关闭，每次锁内检查时间，关闭程序迟到也拒绝超期新请求}{fig:ch12-lock-cutoff}

复制状态机采用相同的资格语义，但时间输入由有权排序命令的准入者取得，再作为命令的一部分复制。跟随者应用命令时读取自己的本机时钟，会让同一日志产生不同结果。各副本使用同一受信判期依据，确定性地核对固定 $t_{\rm deadline}$ 和 $\Gamma$。

排序者在同一串行入口中取得时间区间并分配命令顺序，保存单调的判期进度。它只为 $t_{\mathrm{high}}<t_{\rm deadline}$ 的请求产生可受理命令，已经进入过期边界以后拒绝新的准入。命令应用时再检查 $\Gamma$：关闭命令已经在前，就拒绝其后的旧会话请求。新领导者先恢复已提交前缀和时间进度，再取得可信时间，恢复准入。旧领导者和客户端无法以旧时间戳自行补发一张资格票据。

\CTwelveFigure{ch12-replicated-cutoff}{有权排序者把判期依据和位置共同写入日志，各副本按同一关闭次序应用}{fig:ch12-replicated-cutoff}

因此一条截止前已经排定、后来才提交的命令，可以属于关闭以前的在途受理。一个截止以后才来到排序入口的新请求，即使关闭日志还未提交，也得不到新的准入位置。复制日志保存的是有依据的时序，UTC 截止的真实性来自上述时间前提。若只把 \code{CLOSE(E1)} 放进一个迟到的定时任务，却让此前所有来件一律按 \code{active} 执行，就会留下截止到关闭之间的漏洞~\cite{ongaro}。\footnote{Ongaro 写道“session expiry must be deterministic”。其 LogCabin 示例把领导者时间随命令复制，并让活跃客户端保活。本书采用固定截止，另行明确可信时间及准入次序，E1 关闭后不续期。}
\FloatBarrier
